\ifdefined\gkver\else\def\gkver{student}\fi
\documentclass[\gkver]{gaokaozhenti}
\begin{document}

\chapter{2020年江苏}

\begin{enumerate}
\item
%% number: 1
%% paperName: 2020年江苏高考第 1 题 3 分
%% typeId: 1
%% type: 单选题
%% chapter: 功和能
%% point: 功率
%% method: 无
%% score: 3
%% degree: 950
%% duplicateId: 0
%% body:
质量为 $1.5\times10^{3}\Ukg$ 的汽车在水平路面上匀速行驶，速度为 $20\Ums$，受到的阻力大小为 $1.8\times10^{3}\UN$。此时，汽车发动机输出的实际功率是\xzanswer{C}

\fourchoices[answer=C]
{$90\UW$}
{$30\UkW$}
{$36\UkW$}
{$300\UkW$}

%% answer: C
%% memo:
\memoanswer{
汽车匀速行驶，牵引力与阻力平衡，$F=f=1.8\times10^{3}\UN$，汽车发动机输出的实际功率 $P=Fv=1.8\times10^{3}\times20\UW=36\UkW$，故 C 正确，ABD 错误。
}

\item
%% number: 2
%% paperName: 2020年江苏高考第 2 题 3 分
%% typeId: 1
%% type: 单选题
%% chapter: 交流电
%% point: 变压器
%% method: 无
%% score: 3
%% degree: 850
%% duplicateId: 0
%% body:
电流互感器是一种测量电路中电流的变压器，工作原理如图所示。其原线圈匝数较少，串联在电路中，副线圈匝数较多，两端接在电流表上。则电流互感器\xzanswer{D}

\onepicture[w=3.00cm]{figs/02.png}

\fourchoices[answer=D]
{是一种降压变压器}
{能测量直流电路的电流}
{原、副线圈电流的频率不同}
{副线圈的电流小于原线圈的电流}

%% answer: D
%% memo:
\memoanswer{
A．原线圈匝数较少，电流互感器是一种降电流的变压器，A 选项错误。

B．电流互感器的工作原理是电磁感应中的互感现象，只可以测量交变电流，B 选项错误。

C．电流互感器不会改变电流的频率，只改变电流，故原、副线圈电流的频率相同，C 选项错误。

D．原线圈匝数较少，电流互感器是一种降电流的变压器，副线圈的电流小于原线圈的电流，D 正确。

故选 D。
}

\item
%% number: 3
%% paperName: 2020年江苏高考第 3 题 3 分
%% typeId: 1
%% type: 单选题
%% chapter: 电磁感应
%% point: 楞次定律
%% method: 无
%% score: 3
%% degree: 650
%% duplicateId: 0
%% body:
如图所示，两匀强磁场的磁感应强度 $B_{1}$ 和 $B_{2}$ 大小相等、方向相反。金属圆环的直径与两磁场的边界重合。下列变化会在环中产生顺时针方向感应电流的是\xzanswer{B}

\onepicture[w=4.00cm]{\includesvg{figs/03}}

\fourchoices[answer=B]
{同时增大 $B_{1}$ 减小 $B_{2}$}
{同时减小 $B_{1}$ 增大 $B_{2}$}
{同时以相同的变化率增大 $B_{1}$ 和 $B_{2}$}
{同时以相同的变化率减小 $B_{1}$ 和 $B_{2}$}

%% answer: B
%% memo:
\memoanswer{
AB．产生顺时针方向的感应电流则感应磁场的方向垂直纸面向里。由楞次定律可知，圆环中的净磁通量变化为向里磁通量减少或者向外的磁通量增多，A 错误，B 正确。

CD．同时以相同的变化率增大 $B_{1}$ 和 $B_{2}$，或同时以相同的变化率减小 $B_{1}$ 和 $B_{2}$，两个磁场的磁通量总保持大小相同，所以总磁通量为 0，不会产生感应电流，CD 错误。

故选 B。
}

\item
%% number: 4
%% paperName: 2020年江苏高考第 4 题 3 分
%% typeId: 1
%% type: 单选题
%% chapter: 功和能
%% point: 功能中的图象
%% method: 无
%% score: 3
%% degree: 550
%% duplicateId: 0
%% body:
如图所示，一小物块由静止开始沿斜面向下滑动，最后停在水平地面上。斜面和地面平滑连接，且物块与斜面、物块与地面间的动摩擦因数均为常数。该过程中，物块的动能 $E_{k}$ 与水平位移 $x$ 关系的图象是\xzanswer{A}

\onepicture[w=3.00cm, align=right]{\includesvg{figs/04a}}

\fourchoices[answer=A, ispicture=true, h=2.20cm]
{figs/04b.png}
{figs/04c.png}
{figs/04d.png}
{figs/04e.png}

%% answer: A
%% memo:
\memoanswer{
由动能定理可知 $E_{k}-x$ 图象的斜率表示物块受到的合外力，因为物块与斜面、物块与地面间的动摩擦因数均为常数，所以物块在斜面、地面上运动时受到的合外力均为恒力，$E_{k}-x$ 图象为两段直线，故 A 正确，BCD 错误。

故选 A。
}

\item
%% number: 5
%% paperName: 2020年江苏高考第 5 题 3 分
%% typeId: 1
%% type: 单选题
%% chapter: 运动和力
%% point: 连接体
%% method: 无
%% score: 3
%% degree: 550
%% duplicateId: 0
%% body:
中欧班列在欧亚大陆开辟了“生命之路”，为国际抗疫贡献了中国力量。某运送防疫物资的班列由 40 节质量相等的车厢组成，在车头牵引下，列车沿平直轨道匀加速行驶时，第 2 节对第 3 节车厢的牵引力为 $F$。若每节车厢所受摩擦力、空气阻力均相等，则倒数第 3 节对倒数第 2 节车厢的牵引力为\xzanswer{C}

\fourchoices[answer=C]
{$F$}
{$\frac{19F}{20}$}
{$\frac{F}{19}$}
{$\frac{F}{20}$}

%% answer: C
%% memo:
\memoanswer{
根据题意可知第 2 节车厢对第 3 节车厢的牵引力为 $F$，因为每节车厢质量相等，阻力相同，故对第 2 节车厢以后的 38 节车厢，根据牛顿第二定律有 $F-38f=38ma$；设倒数第 3 节车厢对倒数第 2 节车厢的牵引力为 $F_{1}$，则对最后 2 节车厢有 $F_{1}-2f=2ma$；联立解得 $F_{1}=\frac{F}{19}$。

故选 C。
}

\item
%% number: 6
%% paperName: 2020年江苏高考第 6 题 4 分
%% typeId: 2
%% type: 多选题
%% chapter: 电路
%% point: 动态分析
%% method: 无
%% score: 4
%% degree: 650
%% duplicateId: 0
%% body:
某汽车的电源与启动电机、车灯连接的简化电路如图所示。当汽车启动时，开关 S 闭合，电机工作，车灯突然变暗，此时\xzanswer{ABD}

\onepicture[w=4.20cm]{\includesvg{figs/06}}

\fourchoices[answer=ABD]
{车灯的电流变小}
{路端电压变小}
{电路的总电流变小}
{电源的总功率变大}

%% answer: ABD
%% memo:
\memoanswer{
A．开关闭合时，车灯变暗，故流过车灯的电流 $I_{\text{灯}}$ 变小，A 正确；

B．电路的路端电压 $U_{\text{路}}=U_{\text{灯}}=I_{\text{灯}}R_{\text{灯}}$，$I_{\text{灯}}$ 变小，路端电压变小，B 正确；

C．总电流即干路电流 $I_{\text{干}}=\frac{U_{\text{内}}}{r}=\frac{E-U_{\text{路}}}{r}$，$U_{\text{路}}$ 减小，干路电流增大，C 错误；

D．电源总功率 $P_{\text{总}}=EI_{\text{干}}$，$I_{\text{干}}$ 增大，总功率变大，D 正确。

故选 ABD。
}

\item
%% number: 7
%% paperName: 2020年江苏高考第 7 题 4 分
%% typeId: 2
%% type: 多选题
%% chapter: 曲线运动
%% point: 人造地球卫星
%% method: 无
%% score: 4
%% degree: 550
%% duplicateId: 0
%% body:
甲、乙两颗人造卫星质量相等，均绕地球做圆周运动，甲的轨道半径是乙的 2 倍。下列应用公式进行的推论正确的有\xzanswer{CD}

\fourchoices[answer=CD]
{由 $v=\sqrt{gR}$ 可知，甲的速度是乙的 $\sqrt{2}$ 倍}
{由 $a=\omega^{2}r$ 可知，甲的向心加速度是乙的 2 倍}
{由 $F=G\frac{Mm}{r^{2}}$ 可知，甲的向心力是乙的 $\frac{1}{4}$}
{由 $\frac{r^{3}}{T^{2}}=k$ 可知，甲的周期是乙的 $2\sqrt{2}$ 倍}

%% answer: CD
%% memo:
\memoanswer{
卫星绕地球做圆周运动，万有引力提供向心力，则 $F_{\text{向}}=\frac{GMm}{r^{2}}=\frac{mv^{2}}{r}=m\omega^{2}r=m\frac{4\pi^{2}}{T^{2}}r=ma$。

A．因为在不同轨道上 $g$ 是不一样的，故不能根据 $v=\sqrt{gR}$ 得出甲、乙速度的关系，卫星的运行线速度 $v=\sqrt{\frac{GM}{r}}$，代入数据可得 $\frac{v_{\text{甲}}}{v_{\text{乙}}}=\sqrt{\frac{r_{\text{乙}}}{r_{\text{甲}}}}=\frac{\sqrt{2}}{2}$，故 A 错误；

B．因为在不同轨道上两卫星的角速度不一样，故不能根据 $a=\omega^{2}r$ 得出两卫星加速度的关系，卫星的运行加速度 $a=\frac{GM}{r^{2}}$，代入数据可得 $\frac{a_{\text{甲}}}{a_{\text{乙}}}=\frac{r_{\text{乙}}^{2}}{r_{\text{甲}}^{2}}=\frac{1}{4}$，故 B 错误；

C．根据 $F_{\text{向}}=G\frac{Mm}{r^{2}}$，两颗人造卫星质量相等，可得 $\frac{F_{\text{向甲}}}{F_{\text{向乙}}}=\frac{r_{\text{乙}}^{2}}{r_{\text{甲}}^{2}}=\frac{1}{4}$，故 C 正确；

D．两卫星均绕地球做圆周运动，根据开普勒第三定律 $\frac{r^{3}}{T^{2}}=k$，可得 $\frac{T_{\text{甲}}}{T_{\text{乙}}}=\sqrt{\frac{r_{\text{甲}}^{3}}{r_{\text{乙}}^{3}}}=2\sqrt{2}$，故 D 正确。

故选 CD。
}

\item
%% number: 8
%% paperName: 2020年江苏高考第 8 题 4 分
%% typeId: 2
%% type: 多选题
%% chapter: 曲线运动
%% point: 平抛运动
%% method: 无
%% score: 4
%% degree: 550
%% duplicateId: 0
%% body:
如图所示，小球 A、B 分别从 2$l$ 和 $l$ 的高度水平抛出后落地，上述过程中 A、B 的水平位移分别为 $l$ 和 2$l$。忽略空气阻力，则\xzanswer{AD}

\onepicture[w=4.20cm]{figs/08.png}

\fourchoices[answer=AD]
{A 和 B 的位移大小相等}
{A 的运动时间是 B 的 2 倍}
{A 的初速度是 B 的 $\frac{1}{2}$}
{A 的末速度比 B 的大}

%% answer: AD
%% memo:
\memoanswer{
小球 A 和 B 的位移大小都为 $\sqrt{(2l)^{2}+l^{2}}=\sqrt{5}l$，故 A 正确；由 $h=\frac{1}{2}gt^{2}$ 可知小球 A 的运动时间是小球 B 的 $\sqrt{2}$ 倍，故 B 错误；由 $h=\frac{1}{2}gt^{2}$、$x=v_{0}t$ 可知小球 A 的初速度是小球 B 的 $\frac{\sqrt{2}}{4}$，故 C 错误；由 $h=\frac{1}{2}gt^{2}$、$x=v_{0}t$、$v_{y}=gt$、$v=\sqrt{v_{0}^{2}+v_{y}^{2}}$ 可知小球 A 的末速度为 $v_{A}=\frac{\sqrt{17}}{2}\sqrt{gl}$，小球 B 的末速度为 $v_{B}=2\sqrt{gl}$，因此小球 A 的末速度比小球 B 的大，故 D 正确。
}

\item
%% number: 9
%% paperName: 2020年江苏高考第 9 题 4 分
%% typeId: 2
%% type: 多选题
%% chapter: 电场
%% point: 带电粒子的圆周运动
%% method: 无
%% score: 4
%% degree: 550
%% duplicateId: 0
%% body:
如图所示，绝缘轻杆的两端固定带有等量异号电荷的小球（不计重力）。开始时，两小球分别静止在 A、B 位置。现外加一匀强电场 $E$，在静电力作用下，小球绕轻杆中点 O 转到水平位置。取 O 点的电势为 0。下列说法正确的有\xzanswer{AB}

\onepicture[w=4.30cm]{figs/09.png}

\fourchoices[answer=AB]
{电场 $E$ 中 A 点电势低于 B 点}
{转动中两小球的电势能始终相等}
{该过程静电力对两小球均做负功}
{该过程两小球的总电势能增加}

%% answer: AB
%% memo:
\memoanswer{
A．沿着电场线方向，电势降低，A 正确；

B．由于 O 点的电势为 0，根据匀强电场的对称性 $\varphi_{A}=-\varphi_{B}$，又 $q_{A}=-q_{B}$，$E_{p}=q\varphi$，所以 $E_{pA}=E_{pB}$，B 正确；

CD．A、B 位置的小球受到的静电力分别水平向右、水平向左，绝缘轻杆逆时针旋转，两小球静电力对两小球均做正功，电场力做正功，电势能减少，CD 错误。

故选 AB。
}

\item
%% number: 10
%% paperName: 2020年江苏高考第 10 题 8 分
%% typeId: 4
%% type: 实验题
%% chapter: 电路
%% point: 伏安法测电阻
%% method: 无
%% score: 8
%% degree: 550
%% duplicateId: 0
%% body:
某同学描绘一种电子元件的 $I-U$ 关系图象，采用的实验电路图如题图甲所示，V 为电压表，mA 为电流表，$E$ 为电源（电动势约 $6\UV$），$R$ 为滑动变阻器（最大阻值 $20\UO$），$R_{0}$ 为定值电阻，S 为开关。

\twopicture[numstyle=chinese, wA=4.20cm, wB=4.60cm]{figs/10a.png}{figs/10b.png}

\begin{enumerate}
\item
请用笔画线代替导线，将题图乙所示的实物电路连接完整\tkanswer[3cm]{}。
\item
调节滑动变阻器，记录电压表和电流表的示数如下表：

\begin{center}
\begin{tabular}{|c|c|c|c|c|c|c|c|}
\hline
电压 $U/\UV$ & 0.000 & 0.250 & 0.500 & 0.650 & 0.700 & 0.725 & 0.750 \\
\hline
电流 $I/\UmA$ & 0.00 & 0.10 & 0.25 & 0.60 & 1.70 & 4.30 & 7.50 \\
\hline
\end{tabular}
\end{center}

请根据表中的数据，在方格纸上作出该元件的 $I-U$ 图线\tkanswer[3cm]{}。

\onepicture[w=7.00cm]{figs/10c.png}
\item
根据作出的 $I-U$ 图线可知，该元件是\tkanswer{非线性}（选填“线性”或“非线性”）元件。
\item
在上述测量中，如果用导线代替电路中的定值电阻 $R_{0}$，会导致的两个后果是\tkanswer{BC}。

A．电压和电流的测量误差增大

B．可能因电流过大烧坏待测元件

C．滑动变阻器允许的调节范围变小

D．待测元件两端电压的可调节范围变小
\end{enumerate}

%% answer:
\jdanswer{
\begin{enumerate}
\item
实物连线如答图甲所示。
\item
$I-U$ 图线如答图乙所示。
\item
非线性。
\item
BC。
\end{enumerate}
}

%% memo:
\memoanswer{
（1）根据电路图连接实物图，画线不能跨过仪器、不能交叉，滑动变阻器采用分压式接法，实物图如答图甲所示。

\onepicture[num=甲, labelprefix={答图}, w=4.50cm]{figs/10_an1.png}

（2）根据表格数据描点连线，让尽可能多的点在线上，线尽可能平滑地连接各点，偏离过大的点应舍去，则该元件的 $I-U$ 图线如答图乙所示。

\onepicture[num=乙, labelprefix={答图}, w=6.00cm]{figs/10_an2.png}

（3）根据作出的 $I-U$ 图线可知该元件是非线性元件。

（4）定值电阻 $R_{0}$ 在电路中起到分压、保护电路的作用，如果用导线代替 $R_{0}$，电压和电流的测量误差不会增大，待测元件两端电压的可调节范围不变，会导致滑动变阻器允许的调节范围变小，也可能因电流过大烧坏待测元件，故 AD 错误，BC 正确。
}

\item
%% number: 11
%% paperName: 2020年江苏高考第 11 题 10 分
%% typeId: 4
%% type: 实验题
%% chapter: 直线运动
%% point: 自由落体运动
%% method: 无
%% score: 10
%% degree: 550
%% duplicateId: 0
%% body:
疫情期间“停课不停学”，小明同学在家自主开展实验探究。用手机拍摄物体自由下落的视频，得到分帧图片，利用图片中小球的位置来测量当地的重力加速度，实验装置如题图甲所示。

\onepicture[num=甲, w=3.40cm]{figs/11a.png}

\begin{enumerate}
\item
家中有乒乓球、小塑料球和小钢球，其中最适合用作实验中下落物体的是\tkanswer{小钢球}。
\item
下列主要操作步骤的正确顺序是\tkanswer{①③④②}。（填写各步骤前的序号）

①把刻度尺竖直固定在墙上

②捏住小球，从刻度尺旁静止释放

③手机固定在三角架上，调整好手机镜头的位置

④打开手机摄像功能，开始摄像
\item
停止摄像，从视频中截取三帧图片，图片中的小球和刻度如题图乙所示。已知所截取的图片相邻两帧之间的时间间隔为 $\frac{1}{6}\Us$，刻度尺的分度值是 $1\Umm$，由此测得重力加速度为\tkanswer{9.61} $\Umsq$。

\onepicture[num=乙, w=7.50cm]{figs/11b.png}
\item
在某次实验中，小明释放小球时手稍有晃动，视频显示小球下落时偏离了竖直方向。从该视频中截取图片，\tkanswer{仍能}（选填“仍能”或“不能”）用（3）问中的方法测出重力加速度。
\end{enumerate}

%% answer:
\jdanswer{
\begin{enumerate}
\item
小钢球
\item
①③④②
\item
9.61（9.5～9.7 均可）
\item
仍能
\end{enumerate}
}

%% memo:
\memoanswer{
（1）要测量当地重力加速度需要尽量减小空气阻力的影响，所以密度大、体积小的小钢球最适合。

（2）要完成实验首先应该将刻度尺竖直固定在墙上，安装好三脚架，调整好手机摄像头的位置；因为下落时间很短，所以一定要先打开摄像头开始摄像，然后再将小球从刻度尺旁静止释放，故正确的操作顺序为固定刻度尺、调整手机镜头、开始摄像、释放小球。

（3）由三张图片读出小球所在位置的刻度分别为 $2.50\Ucm$，$26.50\Ucm$，$77.20\Ucm$；小球做自由落体运动，根据 $\Delta x=gT^{2}$ 可得
$g=\frac{\Delta x}{T^{2}}=\frac{(77.20-26.50)\times10^{-2}-(26.50-2.50)\times10^{-2}}{\left(\frac{1}{6}\right)^{2}}\Umsq=9.61\Umsq$。

（4）因为就算小球偏离了竖直方向，但是小球在竖直方向上的运动依然是自由落体运动，对实验结果无影响，故仍能用前面的方法测量出重力加速度。
}

\item
%% number: 12
%% paperName: 2020年江苏高考第 12 题 4 分
%% typeId: 6
%% type: 计算题
%% chapter: 运动和力
%% point: 动量守恒定律
%% method: 无
%% score: 4
%% degree: 850
%% duplicateId: 0
%% body:
本题为选修 3-5 模块。

\begin{enumerate}
\item
“测温枪”（学名“红外线辐射测温仪”）具有响应快、非接触和操作方便等优点。它是根据黑体辐射规律设计出来的，能将接收到的人体热辐射转换成温度显示。若人体温度升高，则人体热辐射强度 $I$ 及其极大值对应的波长 $\lambda$ 的变化情况是\xzanswer{B}

\fourchoices[answer=B]
{$I$ 增大，$\lambda$ 增大}
{$I$ 增大，$\lambda$ 减小}
{$I$ 减小，$\lambda$ 增大}
{$I$ 减小，$\lambda$ 减小}
\item
大量处于某激发态的氢原子辐射出多条谱线，其中最长和最短波长分别为 $\lambda_{1}$ 和 $\lambda_{2}$，则该激发态与基态的能量差为\tkanswer{$\frac{hc}{\lambda_{2}}$}，波长为 $\lambda_{1}$ 的光子的动量为\tkanswer{$\frac{h}{\lambda_{1}}$}。（已知普朗克常量为 $h$，光速为 $c$）
\item
一只质量为 $1.4\Ukg$ 的乌贼吸入 $0.1\Ukg$ 的水，静止在水中。遇到危险时，它在极短时间内把吸入的水向后全部喷出，以 $2\Ums$ 的速度向前逃窜。求该乌贼喷出的水的速度大小 $v$。
\end{enumerate}

%% answer:
\jdanswer{
\begin{enumerate}
\item
B
\item
$\frac{hc}{\lambda_{2}}$，$\frac{h}{\lambda_{1}}$
\item
$28\Ums$
\end{enumerate}
}

%% memo:
\memoanswer{
（1）由黑体辐射规律可知，若人体温度升高，人体热辐射强度 $I$ 增大，其极大值对应的波长 $\lambda$ 减小，故 B 正确，ACD 错误。

（2）根据 $c=\lambda\nu$ 可知波长越短，对应光子的频率越大，对应跃迁的能级差越大；可知最短波长 $\lambda_{2}$ 对应基态到激发态的能量差最大，结合 $E=h\nu$ 得 $\Delta E=h\nu_{2}=\frac{hc}{\lambda_{2}}$；波长为 $\lambda_{1}$ 对应的光子动量为 $p_{1}=\frac{h}{\lambda_{1}}$。

（3）乌贼喷水过程时间较短，内力远大于外力；选取乌贼逃窜的方向为正方向，根据动量守恒定律得 $0=Mv_{1}-mv_{2}$，解得喷出水的速度大小为 $v_{2}=\frac{Mv_{1}}{m}=\frac{1.4\times2}{0.1}\Ums=28\Ums$。
}

\item
%% number: 13
%% paperName: 2020年江苏高考第 13 题 4 分
%% typeId: 6
%% type: 计算题
%% chapter: 光学
%% point: 光的电磁说
%% method: 无
%% score: 4
%% degree: 850
%% duplicateId: 0
%% body:
本题包括 A、B 两小题，请选定其中一小题，并在相应的答题区域内作答。若多做，则按 A 小题评分。

\textbf{A．}（选修 3-3）（12 分）

\begin{enumerate}
\item
玻璃的出现和使用在人类生活里已有四千多年的历史，它是一种非晶体。下列关于玻璃的说法正确的有\xzanswer{AC}

\fourchoices[answer=AC]
{没有固定的熔点}
{天然具有规则的几何形状}
{沿不同方向的导热性能相同}
{分子在空间上周期性排列}
\item
一瓶酒精用了一些后，把瓶盖拧紧，不久瓶内液面上方形成了酒精的饱和汽，此时\tkanswer{有}（选填“有”或“没有”）酒精分子从液面飞出。当温度升高时，瓶中酒精饱和汽的密度\tkanswer{增大}（选填“增大”“减小”或“不变”）。
\item
一定质量的理想气体从状态 A 经状态 B 变化到状态 C，其 $p-\frac{1}{V}$ 图象如图所示，求该过程中气体吸收的热量 $Q$。

\onepicture[w=4.20cm, align=right]{\includesvg{figs/13}}
\end{enumerate}

\textbf{B．}（选修 3-4）（12 分）

\begin{enumerate}
\item
电磁波广泛应用在现代医疗中。下列属于电磁波应用的医用器械有\xzanswer{AB}

\fourchoices[answer=AB]
{杀菌用的紫外灯}
{拍胸片的 X 光机}
{治疗咽喉炎的超声波雾化器}
{检查血流情况的“彩超”机}
\item
我国的光纤通信技术处于世界领先水平。光纤内芯（内层玻璃）的折射率比外套（外层玻璃）的\tkanswer{大}（选填“大”或“小”）。某种光纤的内芯在空气中全反射的临界角为 $43^{\circ}$，则该内芯的折射率为\tkanswer{1.5}。（取 $\sin43^{\circ}=0.68$，$\cos43^{\circ}=0.73$，结果保留 2 位有效数字）
\item
国际宇航联合会将 2020 年度“世界航天奖”授予我国“嫦娥四号”任务团队。“嫦娥四号”任务创造了多项世界第一。在探月任务中，“玉兔二号”月球车朝正下方发射一束频率为 $f$ 的电磁波，该电磁波分别在月壤层的上、下表面被反射回来，反射波回到“玉兔二号”的时间差为 $\Delta t$。已知电磁波在月壤层中传播的波长为 $\lambda$，求该月壤层的厚度 $d$。
\end{enumerate}

%% answer:
\jdanswer{
\begin{enumerate}
\item
A：（1）AC；（2）有，增大；（3）$Q=2\times10^{5}\UJ$。
\item
B：（1）AB；（2）大，1.5；（3）$d=\frac{f\lambda\Delta t}{2}$。
\end{enumerate}
}

%% memo:
\memoanswer{
\textbf{A．}（1）玻璃是非晶体，没有固定的熔点，不具有规则的几何形状，分子在空间上没有周期性排列，物理性质具有各向同性，即沿不同方向的导热性能相同，故 AC 正确，BD 错误。

（2）与液体处于动态平衡的蒸汽叫作饱和汽，所以此时有酒精分子从液面飞出；当温度升高时，饱和汽中酒精分子数增多，体积不变，所以瓶中酒精饱和汽的密度增大。

（3）根据 $p-\frac{1}{V}$ 图象可知状态 A 和状态 C 温度相同，内能相同；故从 A 经 B 到 C 过程中气体吸收的热量等于气体对外所做的功。根据图象可知状态 A 到状态 B 为等压过程，气体对外做功为 $W_{1}=p\Delta V=2\times10^{5}\times(2-1)\UJ=2\times10^{5}\UJ$，状态 B 到状态 C 为等容变化，气体不做功；故 A 经 B 到 C 过程中气体吸收的热量为 $Q=W_{1}=2\times10^{5}\UJ$。

\textbf{B．}（1）杀菌用的紫外灯用的是紫外线，拍胸片的 X 光机用的是 X 射线，都是电磁波的应用，故 AB 正确；超声波雾化器、“彩超”机都是超声波的应用，都是机械波的应用，故 CD 错误。

（2）光纤通信利用了光的全反射现象，光由光密介质射入光疏介质时才会发生全反射，所以光纤内芯的折射率比外套的大；该内芯的折射率为 $n=\frac{1}{\sin C}=\frac{1}{\sin43^{\circ}}=\frac{1}{0.68}\approx1.5$。

（3）电磁波在土壤中的传播速度 $v=f\lambda$，由题意得 $2d=v\Delta t$，解得 $d=\frac{f\lambda\Delta t}{2}$。
}

\item
%% number: 14
%% paperName: 2020年江苏高考第 14 题 15 分
%% typeId: 6
%% type: 计算题
%% chapter: 电磁感应
%% point: 线圈过有界磁场
%% method: 无
%% score: 15
%% degree: 650
%% duplicateId: 0
%% body:
如图所示，电阻为 $0.1\UO$ 的正方形单匝线圈 $abcd$ 的边长为 $0.2\Um$，$bc$ 边与匀强磁场边缘重合。磁场的宽度等于线圈的边长，磁感应强度大小为 $0.5\UT$。在水平拉力作用下，线圈以 $8\Ums$ 的速度向右穿过磁场区域。求线圈在上述过程中：

\begin{enumerate}
\item
感应电动势的大小 $E$；
\item
所受拉力的大小 $F$；
\item
感应电流产生的热量 $Q$。
\end{enumerate}

\onepicture[w=5.20cm, align=right]{\includesvg{figs/14}}

%% answer:
\jdanswer{
\begin{enumerate}
\item
$0.8\UV$
\item
$0.8\UN$
\item
$0.32\UJ$
\end{enumerate}
}

%% memo:
\memoanswer{
（1）由题意可知，当线框切割磁感线时产生的电动势为 $E=BLv=0.5\times0.2\times8\UV=0.8\UV$。

（2）因为线框匀速运动，故所受拉力等于安培力，有 $F=F_{\text{安}}=BIL$，根据闭合电路欧姆定律有 $I=\frac{E}{R}$，结合（1）联立各式代入数据可得 $F=0.8\UN$。

（3）线框穿过磁场所用的时间为 $t=\frac{2L}{v}=\frac{2\times0.2}{8}\Us=0.05\Us$，故线框穿越过程产生的热量为 $Q=I^{2}Rt=\frac{E^{2}}{R}t=\frac{0.8^{2}}{0.1}\times0.05\UJ=0.32\UJ$。
}

\item
%% number: 15
%% paperName: 2020年江苏高考第 15 题 16 分
%% typeId: 6
%% type: 计算题
%% chapter: 功和能
%% point: 整体机械能守恒
%% method: 无
%% score: 16
%% degree: 450
%% duplicateId: 0
%% body:
如图所示，鼓形轮的半径为 $R$，可绕固定的光滑水平轴 O 转动。在轮上沿相互垂直的直径方向固定四根直杆，杆上分别固定有质量为 $m$ 的小球，球与 O 的距离均为 2$R$。在轮上绕有长绳，绳上悬挂着质量为 $M$ 的重物。重物由静止下落，带动鼓形轮转动。重物落地后鼓形轮匀速转动，转动的角速度为 $\omega$。绳与轮之间无相对滑动，忽略鼓形轮、直杆和长绳的质量，不计空气阻力，重力加速度为 $g$。求：

\begin{enumerate}
\item
重物落地后，小球线速度的大小 $v$；
\item
重物落地后一小球转到水平位置 A，此时该球受到杆的作用力的大小 $F$；
\item
重物下落的高度 $h$。
\end{enumerate}

\onepicture[w=4.20cm, align=right]{figs/15.png}

%% answer:
\jdanswer{
\begin{enumerate}
\item
$v=2\omega R$
\item
$F=\sqrt{(2m\omega^{2}R)^{2}+(mg)^{2}}$
\item
$H=\frac{(M+16m)R^{2}\omega^{2}}{2Mg}$
\end{enumerate}
}

%% memo:
\memoanswer{
（1）由题意可知当重物落地后鼓形轮转动的角速度为 $\omega$，则根据线速度与角速度的关系 $v=\omega r$，可知小球的线速度为 $v=2\omega R$。

（2）小球匀速转动，当在水平位置时设杆对球的作用力为 $F$，合力提供向心力，则有 $\sqrt{F^{2}-(mg)^{2}}=m\frac{v^{2}}{2R}$，结合（1）可解得杆对球的作用力大小为 $F=\sqrt{(2m\omega^{2}R)^{2}+(mg)^{2}}$。

（3）设重物下落高度为 $H$，重物下落过程中对重物、鼓形轮和小球组成的系统，根据系统机械能守恒可知 $MgH=\frac{1}{2}Mv_{1}^{2}+\frac{1}{2}\cdot4mv^{2}$，而重物的速度等于鼓形轮的线速度，有 $v_{1}=\omega R$，联立各式解得 $H=\frac{(M+16m)R^{2}\omega^{2}}{2Mg}$。
}

\item
%% number: 16
%% paperName: 2020年江苏高考第 16 题 16 分
%% typeId: 6
%% type: 计算题
%% chapter: 磁场
%% point: 带电粒子在磁场中的运动
%% method: 无
%% score: 16
%% degree: 350
%% duplicateId: 0
%% body:
空间存在两个垂直于 $Oxy$ 平面的匀强磁场，$y$ 轴为两磁场的边界，磁感应强度分别为 2$B_{0}$、3$B_{0}$。甲、乙两种比荷不同的粒子同时从原点 O 沿 $x$ 轴正向射入磁场，速度均为 $v$。甲第 1 次、第 2 次经过 $y$ 轴的位置分别为 P、Q，其轨迹如图所示。甲经过 Q 时，乙也恰好同时经过该点。已知甲的质量为 $m$，电荷量为 $q$。不考虑粒子间的相互作用和重力影响。求：

\begin{enumerate}
\item
Q 到 O 的距离 $d$；
\item
甲两次经过 P 点的时间间隔 $\Delta t$；
\item
乙的比荷可能的最小值。
\end{enumerate}

\onepicture[w=4.50cm, align=right]{\includesvg{figs/16}}

%% answer:
\jdanswer{
\begin{enumerate}
\item
$d=\frac{mv}{3qB_{0}}$
\item
$\Delta t=\frac{2\pi m}{qB_{0}}$
\item
$\frac{2q}{m}$
\end{enumerate}
}

%% memo:
\memoanswer{
（1）甲粒子先后在两磁场中做匀速圆周运动，设半径分别为 $r_{1}$、$r_{2}$，由半径 $r=\frac{mv}{qB}$ 得 $r_{1}=\frac{mv}{2qB_{0}}$，$r_{2}=\frac{mv}{3qB_{0}}$，且 $d=2r_{1}-2r_{2}$，联立解得 $d=\frac{mv}{3qB_{0}}$。

（2）甲粒子先后在两磁场中做匀速圆周运动，设运动时间分别为 $t_{1}$、$t_{2}$，由 $T=\frac{2\pi m}{qB}$ 得 $t_{1}=\frac{\pi m}{2qB_{0}}$，$t_{2}=\frac{\pi m}{3qB_{0}}$，且 $\Delta t=2t_{1}+3t_{2}$，联立解得 $\Delta t=\frac{2\pi m}{qB_{0}}$。

（3）乙粒子周期性地先后在两磁场中做匀速圆周运动，若经过两磁场的次数均为 $n$（$n=1,2,3,\cdots$），相遇时，有 $n\frac{m'v}{3q'B_{0}}=d$，$n\frac{5\pi m'}{6q'B_{0}}=t_{1}+t_{2}$，解得 $\frac{q'}{m'}=n\frac{q}{m}$。根据题意，$n=1$ 舍去；当 $n=2$ 时，$\frac{q'}{m'}$ 有最小值 $\left(\frac{q'}{m'}\right)_{\min}=\frac{2q}{m}$。若先后经过右侧、左侧磁场的次数分别为 $(n+1)$、$n$（$n=1,2,3,\cdots$），经分析不可能相遇。综上分析，比荷的最小值为 $\frac{2q}{m}$。
}

\end{enumerate}

\end{document}
